\subsection{Расстояние между прямой и плоскостью}

Найдите расстяние между плоскостью
\[
	\pi =
	a + \left<w_1, w_2\right>
	=
	\ee
	\Vect{4 \\ 1 \\ -2 \\ 5}
	+
	\left<
		\ee
		\Vect{
			0 \\ 1 \\ -2 \\ 4
		},
		\ 
		\ee
		\Vect{
			1 \\ 1 \\ -3 \\ 3
		}
	\right>
	,
\]
где угловые скобки означают линейную оболочку, и прямой
\[
	\ell =
	b +
	\left<
		v
	\right>
	=
	\ee
	\Vect{2 \\ 3 \\ 2 \\ 9}
	+
	\left<
		\ee
		\Vect{
			0 \\ -1 \\ 2 \\ -5
		}
	\right>
	.
\]

\begin{proofm}
	Снова сначала поговорим о трёхмерном случае: это поможет сориентироваться в многомерии.

	Если нам были заданы две скрещивающиеся прямые~$a + \anspan{v}$ и~$b + \anspan{w}$, как можно было бы поступить? Можно организовать вектор~$b - a$, который гарантированно соединяет по одной точке каждой прямой, а затем параллельным переносом перевести направляющие векторы прямых~$v$ и~$w$ и $b - a$ в одно начало. Получится параллелепипед с образующими~$(v, w, b - a)$. Искомое расстояние между прямыми --- это высота параллелепипеда. Как её найти? Можно поделить объём на площадь основания, образованного~$v$ и~$w$:
	\[
		h =
		\frac{
			\vol(v, w, b-a)
		}{
			\area (v, w)
		}
		.
	\]
	При этом~$\vol(v, w, b-a)$ можно высчитать как смешанное произведение (в координатах это будет просто определитель матрицы, в строках или столбцах которой расположены координаты векторов), а $\area(v, w)$ --- как модуль векторного произведения~$\abs{[v, w]}$.

	Вообще, \emph{определители считают $k$-мерный объём}.

	Но в многомерии нет понятия векторного произведения. Зато есть понятие грассманова произведения. Для этого можно было бы организовать отдельную главу, но я скажу только, как с ним работать (хотя, возможно, найдёт вдохновение и всё же напишу её).

	\subsubsection{Грассманова алгебра}

	Всюду в этом подразделе векторы обозначаются полужирными буквами, а греческими обозначаются числа (элементы поля или кольца).

	Грассманово произведение обозначается~$\wedge$ \noatt{и индуцируется тензорным произведением, факторизованным по соотношениям кососимметричности}. Иными словами, для любых векторов~$\vec{a}$, $\vec{b}$
	\begin{align*}
		\vec{a} \wedge \vec{a} &= 0,
		\\
		(\vec{a} + \vec{b}) \wedge \vec{c} &=
		\vec{a} \wedge \vec{c} +
		\vec{b} \wedge \vec{c},
		\\
		\vec{a} \wedge (\vec{b} + \vec{c}) &=
		\vec{a} \wedge \vec{b} +
		\vec{a} \wedge \vec{c},
		\\
		(\alpha \vec{a}) \wedge \vec{b} &=
		\alpha (\vec{a} \wedge \vec{b})
		,
		\\
		\vec{a} \wedge (\alpha \vec{b}) &=
		\alpha (\vec{a} \wedge \vec{b})
		.
	\end{align*}
	заметим, что из свойства~$\vec{a} \wedge \vec{a} = 0$ вытекает, что~$\vec{v} \wedge \vec{w} = -\vec{w} \wedge \vec{v}$ (достаточно подставить $\vec{a} := \vec{v} + \vec{w}$).

	Одно из основных свойств грассмановых произведений вытекает из кососимметричности: если~$\sigma \in S_k$ --- перестановка, а~$(\vec{a}_1, \dots, \vec{a}_k)$ --- набор векторов из~$n$-мерноно пространства, то
	\[
		\vec{a}_{\sigma(1)}
		\wedge
		\vec{a}_{\sigma(2)}
		\wedge
		\cdots
		\wedge
		\vec{a}_{\sigma(k)}
		=
		\sgn{\sigma} \cdot
		\vec{a}_{1}
		\wedge
		\vec{a}_{2}
		\wedge
		\cdots
		\wedge
		\vec{a}_{k}
		.
	\]
	
	Из этого следует удобное определение определителя.
	\begin{definition}
		Пусть~$\ee = (\vec{e}_1, \dots \vec{e}_n)$ --- базис $n$-мерного пространства~$V$, а $\baa = (\vec{a}_1, \dots \vec{a}_n)$ --- набор из~$n$ векторов. (Обрати внимание, что векторов столько же, какова размерность пространства!)
		
		Если $\baa = \ee A$, то есть~$A$ --- матрица, состоящая из координат векторов~$\baa$ в базисе~$\ee$, то~$\det{A}$ \emph{определяется} как множитель в равенстве
		\[
			\vec{a}_1 \wcw \vec{a}_n
			=
			\det{A}
			\cdot
			\vec{e}_1 \wcw \vec{e}_n
			.
		\]
	\end{definition}

	Если хочется связать со знакомым определением, милости прошу. Пусть~$A = (a^i_j)$, где верхний индекс нумерует строки, а нижний --- столбцы. Таким образом,
	\[
		a_j =
		\Vect{a_j^1 \\ \dots \\ a_j^n}
	\]
	--- координаты вектора~$\vec{a}_j$, составляющие~$j$-ый стробец матрицы~$A$. Тогда
	\begin{align*}
		\vec{a}_1 \wcw \vec{a}_n
		&=
		\left( 
			\sum_{i_1 = 1}^{n}
				a_1^{i_1}
				\vec{e_{i_1}}
		\right)
		\wcw
		\left( 
			\sum_{i_n = 1}^{n}
				a_n^{i_n}
				\vec{e_{i_n}}
		\right)
		\nlinea{=}
		\sum_{i_1 = 1}^{n}
		\cdots
		\sum_{i_n = 1}^{n}
			a_1^{i_1}
			\cdots
			a_n^{i_n}
			\cdot
			\vec{e}_{i_1}
			\wcw
			\vec{e}_{i_n}
		.
	\end{align*}
	Мы вынесли суммы и числовые множители по линейности. Теперь заметим, что если в наборе индексов~$(i_1, \dots, i_n)$ есть одинаковые, то
	\[
		\vec{e}_{i_1} \wcw \vec{e}_{i_n} = 0
		.
	\]
	Поэтому останутся только слагаемые с попарно разными индексами, и сумму можно записать в виде
	\begin{multline*}
		\sum_{\sigma \in S_n}
		a_1^{\sigma(1)}
		\cdots
		a_n^{\sigma(n)}
		\cdot
		\vec{e}_{\sigma(1)}
		\wcw
		\vec{e}_{\sigma(n)}
		\nline{=}
		\sum_{\sigma \in S_n}
		a_1^{\sigma(1)}
		\cdots
		a_n^{\sigma(n)}
		\cdot
		\sgn{\sigma}
		\cdot
		\vec{e}_{1}
		\wcw
		\vec{e}_{n}
		\nline{=}
		\left(
			\sum_{\sigma \in S_n}
				\sgn{\sigma}
				\cdot
				a_1^{\sigma(1)}
				\cdots
				a_n^{\sigma(n)}
		\right)
		\cdot
		\vec{e}_{1}
		\wcw
		\vec{e}_{n}
		\nline{=}
		\det{A}
		\cdot
		\vec{e}_{1}
		\wcw
		\vec{e}_{n}
		.
	\end{multline*}
	Итак, определение определителя через грассманову алгебру совпадает со старым.

	Спрашивается: а зачем это всё? А затем, чтобы можно было проще работать с минорами, которые кодируют площади~$k$-мерных объёмов в~$n$-мерном пространстве. Если, например, 
	\[
		\baa = (\vec{a}_1, \dots, \vec{a}_k) = \ee A
	\]
	--- набор из~$k < n$ векторов в~$n$-мерном пространстве, то определитель~$A$ взять не получится (матрица не квдаратная), зато получится взять грассманово произведение:
	\begin{align*}
		\vec{a}_1 \wcw \vec{a}_k
		&=
		\left( 
			\sum_{i_1 = 1}^{n}
				a_1^{i_1}
				\vec{e}_{i_1}
		\right)
		\wcw
		\left( 
			\sum_{i_k = 1}^{n}
				a_k^{i_k}
				\vec{e}_{i_k}
		\right)
		\nlinea{=}
		\sum_{i_1 = 1}^{n}
		\cdots
		\sum_{i_k = 1}^{n}
			a_1^{i_1}
			\cdots
			a_k^{i_k}
		\cdot
			\vec{e}_{i_1}
			\wcw
			\vec{e}_{i_k}
		\nlinea{=}
		\sum_{(i_1, \dots, i_k) \in [n]^k \setminus \Delta}
			a_1^{i_1}
			\cdots
			a_k^{i_k}
		\cdot
			\vec{e}_{i_1}
			\wcw
			\vec{e}_{i_k}
		,
	\end{align*}
	где суммирование идёт по всем наборам~$(i_1, \dots, i_k) \in [n]$, попарно не совпадающими друг с другом.

	Частный случай этой ситуации, который для нас важен, --- это многомерный аналог векторного произведения, получающийся при~$k = n-1$. В этом случае
	\begin{align*}
		\vec{a}_1 \wcw \vec{a}_{n-1}
		&=
		\left( 
			\sum_{i_1 = 1}^{n}
				a_1^{i_1}
				\vec{e}_{i_1}
		\right)
		\wcw
		\left( 
			\sum_{i_{n-1} = 1}^{n}
				a_{n-1}^{i_{n-1}}
				\vec{e}_{i_{n-1}}
		\right)
		\nlinea{=}
		\sum_{i_1 = 1}^{n}
		\cdots
		\sum_{i_{n-1} = 1}^{n}
			a_1^{i_1}
			\cdots
			a_{n-1}^{i_{n-1}}
		\cdot
			\vec{e}_{i_1}
			\wcw
			\vec{e}_{i_{n-1}}
		\nlinea{=}
		\sum_{(i_1, \dots, i_{n-1}) \in [n] \setminus \Delta}
			a_1^{i_1}
			\cdots
			a_{n-1}^{i_{n-1}}
		\cdot
			\vec{e}_{i_1}
			\wcw
			\vec{e}_{i_{n-1}}
		.
	\end{align*}
	Суммирование ведётся по попарно не повторяющимся наборам чисел~$(i_1, \dots, i_{n-1})$. Но создать такой набор значит выбрать~$n-1$ число, входящее в набор, из~$n$ чисел, или, что то же самое, \emph{выбрать одно число, которое не входит в набор}. Поэтому суммирование можно вести по~\emph{одному} числу~$i$, \emph{не входящему в грассманов моном}.

	Каждое слагаемого
	\[
		a_1^{i_1}
		\cdots
		a_k^{i_{n-1}}
		\cdot
		\vec{e}_{i_1}
		\wcw
		\vec{e}_{i_{n-1}}
	\]
	умножим на \emph{фиктивный} вектор~$\tta^i \te_i$, где $i$ --- единственное число, не входящее в набор~$(i_1, \dots, i_{n-1})$. Тогда
	\begin{multline*}
		\sum_{(i_1, \dots, i_{n-1}) \in [n] \setminus \Delta}
			a_1^{i_1}
			\cdots
			a_{n-1}^{i_{n-1}}
			\ta_i
		\cdot
			\te_i \wedge
			\vec{e}_{i_1}
			\wcw
			\vec{e}_{i_{n-1}}
		\nline{=}
		\sum_{(i_1, \dots, i_{n-1}) \in [n] \setminus \Delta}
			a_1^{i_1}
			\cdots
			a_{n-1}^{i_{n-1}}
			\ta_i
		\cdot
			\sgn(i_1, \cdots, i_{n-1})
		\cdot
			\te_i \wedge
			\vec{e}_1
			\wcw
			\widehat{\vec{e}}_i
			\wcw
			\vec{e}_n
		\nline{=}
		\sum_{(i_1, \dots, i_{n-1}) \in [n] \setminus \Delta}
			a_1^{i_1}
			\cdots
			a_{n-1}^{i_{n-1}}
			\ta_i
		\cdot
			\sgn(i_1, \cdots, i_{n-1})
			\cdot
			(-1)^{i-1}
		\cdot
			\vec{e}_1
			\wcw
			\te_i
			\wcw
			\vec{e}_n
		\nline{=}
		\sum_{i = 1}^n
			(-1)^{i-1}
			\ta_i
			\left( 
				\sum_{(i_1, \dots, i_{n-1}) \in S_{n-1}}
					a_1^{i_1}
					\cdots
					a_{n-1}^{i_{n-1}}
				\cdot
					\sgn(i_1, \cdots, i_{n-1})
			\right)
			\cdot
				\vec{e}_1
				\wcw
				\te_i
				\wcw
				\vec{e}_n
			.
	\end{multline*}
	Выражение в скобках --- это определитель минора матрицы, в которой вычеркнули первую строку и~$i$-ый столбец. А сумма со знаками таких миноров, умноженных на~$\ta_i$, --- это определитель матрицы
	\[
		\begin{pmatrix}
			\ta_1 & \ta_2 & \dots & \ta_n
			\\
			a_1^1 & a_1^2 & \dots & a_1^n
			\\
			\hdotsfor{4}
			\\
			a_{n-1}^1 & a_{n-1}^2 & \dots & a_{n-1}^n
		\end{pmatrix}
		.
	\]
	Обычно в верхней строке такого определителя стоят формальные базисные векторы:
	\[
		[\vec{a}_1, \dots, \vec{a}_{n-1}] :=
		\begin{determ}
			\vec{e}_1 & \vec{e}_2 & \dots & \vec{e}_n
			\\
			a_1^1 & a_1^2 & \dots & a_1^n
			\\
			\hdotsfor{4}
			\\
			a_{n-1}^1 & a_{n-1}^2 & \dots & a_{n-1}^n
		\end{determ}
		.
	\]
	Такой определитель естественно назвать \tdef{векторным произведением} векторов~$\vec{a}_1$, \dots, $\vec{a}_{n-1}$ в~$n$-мерном векторном пространстве~$V$.
	\begin{remark}
		Определитель можно написать и транспонированным:
		\[
			[\vec{a}_1, \dots, \vec{a}_{n-1}] =
			\begin{determ}
				\vec{e}_1 & a_1^1 & a_2^1 & \dots & a_{n-1}^1 \\
				\vec{e}_2 & a_1^2 & a_2^2 & \dots & a_{n-1}^2 \\
				\vec{e}_3 & a_1^3 & a_2^3 & \dots & a_{n-1}^3 \\
				\hdotsfor{5} \\
				\vec{e}_{n-1} & a_1^{n-1} & a_2^{n-1} & \dots & a_{n-1}^{n-1}
			\end{determ}
			.
		\]
		Это соответствует записи координат векторов в столбик.
	\end{remark}



	% Обозначим~$[n] \setminus i$ набор чисел от~$1$ до~$n$, в котором нет числа~$i$, а для грассмановых мономов положим
	% \[
	% 	\vec{e}_{[n] \setminus i} :=
	% 	\vec{e}_1
	% 	\wcw
	% 	\widehat{
	% 		\vec{e}_i
	% 	}
	% 	\wcw
	% 	\vec{e}_n
	% \]
	% (крышка означает, что множитель опущен).
	% Аналогично обозначим
	% \[
	% 	a^{[n] \setminus i} :=
	% 	a_1^{i_1} \cdots a_{n-1}^{i_{n-1}},
	% \]
	% где $(i_1, \dots, i_{n-1}) = [n] \setminus i$.
	% Тогда
	% \begin{align*}
	% 	\vec{a}_1 \wcw \vec{a}_{n-1}
	% 	&=
	% 	\sum_{i = 1}^{n}
	% 	\left( 
	% 		\sum_{
	% 			(i_1, \dots, i_{n-1})
	% 			\in
	% 			[n] \setminus i
	% 		}
	% 			a_1^{i_1}
	% 			\cdots
	% 			a_k^{i_{n-1}}
	% 		\cdot
	% 			\sgn(i_1, \dots, i_n)
	% 			\vec{e}_{[n] \setminus i}
	% 	\right)
	% 	\nlinea{=}
	% 	\sum_{i = 1}^{n}
	% 		\sum_{
	% 			(i_1, \dots, i_{n-1})
	% 			\in
	% 			[n] \setminus i
	% 		}
	% 	\left( 
	% 		a_1^{i_1}
	% 		\cdots
	% 		a_k^{i_{n-1}}
	% 		\cdot
	% 		\sgn(i_1, \dots, i_n)
	% 	\right)
	% 	\vec{e}_{[n] \setminus i}
	% 	.
	% \end{align*}
	% Если теперь приписать к каждому моному
	% $
	% 	a_1^{i_1}
	% 	\cdots
	% 	a_k^{i_{n-1}}
	% $
	% \emph{фиктивный} векторный множитель~$\tvec{e}_i$, где $i \notin (i_1, \dots, i_{n-1})$, то
	% \begin{align*}
	% 	\vec{a}_1 \wcw \vec{a}_{n-1}
	% 	&=
	% 	\sum_{
	% 		(i_1, \dots, i_{n-1})
	% 		\in
	% 		[n]^{n-1}\setminus\Delta
	% 	}
	% 		a_1^{i_1}
	% 		\cdots
	% 		a_k^{i_{n-1}}
	% 		\cdot
	% 		\sgn(i_1, \dots, i_n)




	% \end{align*}
		

	\subsubsection{Возвращаемся к задаче}
		
		С учётом всего вышесказанного можно написать формулу для нахождения расстояния:
		\begin{equation}
			\label{eq:num:plane_line_dist:dist}
			d =
			\frac{
				\abs{[w_1, w_2, v]}
			}{
				\abs{(w_1, w_2, v, b-a)}
			}.
		\end{equation}
		Поплыли считать:
		\[
			[w_1, w_2, v]
			=
			\begin{determ}
				\vec{e}_1 & 
				\vec{e}_2 & 
				\vec{e}_3 & 
				\vec{e}_4 \\
				0 & 1 & -2 & 4 \\
				1 & 1 & -3 & 3 \\
				0 & -1 & 2 & 5
			\end{determ}
			.
		\]
		Раскроем определитель по первой строке. Множители при~$\vec{e}_k$ --- это определители $3 \times 3$ с учётом знака:
		\begin{align*}
			\begin{determ}
				1 & -2 & 4 \\
				1 & -3 & 3 \\
				-1 & 2 & 5
			\end{determ}
			\cdot
			\vec{e}_1
			&=
			-9 \vec{e}_1
			;
			\\
			\begin{determ}
				0 & -2 & 4 \\
				1 & -3 & 3 \\
				0 &  2 & 5
			\end{determ}
			\cdot
			(-\vec{e}_2)
			&=
			-18 \vec{e}_2
			;
			\\
			\begin{determ}
				0 & 1  & 4 \\
				1 & 1  & 3 \\
				0 & -1 & 5
			\end{determ}
			\cdot
			\vec{e}_3
			&=
			9 \vec{e}_3
			;
			\\
			\\
			\begin{determ}
				0 & 1 & -2 \\
				1 & 1 & -3 \\
				0 & -1 & 2
			\end{determ}
			\cdot
			(-\vec{e}_4)
			&=
			0\vec{e}_4
			.
		\end{align*}
		Итак,
		\[
			[w_1, w_2, v] =
			-9\vec{e}_1 - 18 \vec{e}_2 - 9 \vec{e}_3 + 0\vec{e}_4
			=
			\ee
			\cdot 
			-9
			\Vect{1 \\ 2 \\ 1 \\ 0}.
		\]
		Тогда
		\[
			\norm{[w_1, w_2, v]} =
			\norm{
				\ee
				\cdot 
				(-9)
				\Vect{1 \\ 2 \\ 1 \\ 0}
			}
			=
			9 \cdot 
			\norm{
				\ee
				\cdot 
				\Vect{1 \\ 2 \\ 1 \\ 0}
			}
			=
			9 \cdot (1 + 4 + 1)
			=
			36
			.
		\]
		Теперь посчитаем смешанное произведение --- объём четырёхмерного параллелепипеда. Для этого нам нужны координаты вектора~$b-a$:
		\[
			b-a =
			\ee
			\left[ 
				\Vect{2 \\ 3 \\ 2 \\ 9}
				-
				\Vect{4 \\ 1 \\ -2 \\ 5}
			\right]
			=
			\ee
			\Vect{
				-2 \\ 2 \\ 4 \\ 4
			}
			=
			2
			\cdot
			\ee
			\Vect{
				-1 \\ 1 \\ 2 \\ 2
			}
		\]
		Тогда
		\begin{align*}
			\abs{(w_1, w_2, v, b-a)} =
			\begin{determ}
				0 & 1 & -2 & 4 \\
				1 & 1 & -3 & 3 \\
				0 & -1 & 2 & -5 \\
				-2 & 2 & 4 & 4
			\end{determ}
			=
			-6
			.
		\end{align*}
		Поэтому по формуле~\eqref{eq:num:plane_line_dist:dist}
		\[
			d =
			\frac{36}{\abs{-6}}
			=
			6
			.
			\proofmark
		\]
\end{proofm}
